\newpage
\section{Simulation}

Chaque partie a été simulée manuellement sur Logisim dans le but de visualiser les priorités et les rebouclages des systèmes.
\subsection{Partie 1 : registre 4 bits}

Le tableau~\ref{tab:p1_tests} donne, pour chaque flanc montant de
\texttt{clock\_sti}, la fonction testée et la valeur attendue sur \texttt{reg\_obs}.
Les valeurs observées sur le chronogramme (figure~\ref{fig:p1_chrono})
correspondent toutes aux valeurs attendues.

\begin{table}[H]
  \centering
  \small
  \begin{tabular}{c l c}
    \hline
    \textbf{Flanc} & \textbf{Fonction testée} & \textbf{\texttt{reg\_obs} attendu} \\
    \hline
    15 et 25~µs & \textbf{Reset prioritaire} (\texttt{en\_sh} = \texttt{ser\_r} = 1) & \texttt{0000} \\
    35~µs  & Décalage à droite, \texttt{ser\_r\_sti} = 1 & \texttt{1000} \\
    45~µs  & Décalage à droite, \texttt{ser\_r\_sti} = 1 & \texttt{1100} \\
    55~µs  & Décalage à droite, \texttt{ser\_r\_sti} = 1 & \texttt{1110} \\
    65~µs  & Décalage à droite, \texttt{ser\_r\_sti} = 1 & \texttt{1111} \\
    75 et 85~µs & Maintien (\texttt{en\_sh\_sti} = 0)   & \texttt{1111} \\
    95~µs  & Décalage à droite, \texttt{ser\_r\_sti} = 0 & \texttt{0111} \\
    105~µs & Décalage à droite, \texttt{ser\_r\_sti} = 0 & \texttt{0011} \\
    110~µs & \textbf{Reset asynchrone} & \texttt{0000} \\
    \hline
  \end{tabular}
  \caption{Tests de simulation de \texttt{srg4\_sh\_r}}
  \label{tab:p1_tests}
\end{table}

\capture{p1_chrono_annote}{Chronogramme de \texttt{srg4\_sh\_r\_sim}}{fig:p1_chrono}{}

\needspace{12\baselineskip}
\subsection{Partie 2 : registre 8 bits}

Le tableau~\ref{tab:p2_tests} donne, pour chaque flanc montant de
\texttt{clock\_i}, la fonction testée et la valeur attendue sur \texttt{reg\_o}.
Les valeurs observées sur le chronogramme (figures~\ref{fig:p2_chrono_a} et~\ref{fig:p2_chrono_b})
correspondent toutes aux valeurs attendues.

\begin{table}[H]
  \centering
  \begin{tabular}{c l c}
    \hline
    \textbf{Flanc} & \textbf{Fonction testée} & \textbf{\texttt{reg\_o} attendu} \\
    \hline
    55~µs  & Chargement de \texttt{10110010}                    & \texttt{1011 0010} \\
    85~µs  & Maintien (\texttt{en\_sh\_i} = 0)                  & \texttt{1011 0010} \\
    105~µs & Décalage à droite, \texttt{ser\_r\_i} = 0          & \texttt{0101 1001} \\
    120~µs & Décalage à droite, \texttt{ser\_r\_i} = 1          & \texttt{1010 1100} \\
    130~µs & Décalage à gauche, \texttt{ser\_l\_i} = 1          & \texttt{0101 1001} \\
    140~µs & Décalage à gauche, \texttt{ser\_l\_i} = 0          & \texttt{1011 0010} \\
    150~µs & \textbf{Load prioritaire sur le décalage} (\texttt{11110000}) & \texttt{1111 0000} \\
    165~µs & Décalage à droite, \texttt{ser\_r\_i} = 1          & \texttt{1111 1000} \\
    175~µs & Décalage à droite, \texttt{ser\_r\_i} = 1          & \texttt{1111 1100} \\
    185~µs & \textbf{Reset asynchrone}      & \texttt{0000 0000} \\
    \hline
  \end{tabular}
  \caption{Tests de simulation de \texttt{srg8\_ld\_sh}}
  \label{tab:p2_tests}
\end{table}

\capture{p2_chrono_a}{Chronogramme de \texttt{srg8\_ld\_sh\_sim} (1/2) : de 0 à 108~µs}{fig:p2_chrono_a}{}
\capture{p2_chrono_b}{Chronogramme de \texttt{srg8\_ld\_sh\_sim} (2/2) : de 103 à 200~µs}{fig:p2_chrono_b}{}

\needspace{12\baselineskip}
\subsection{Partie 3 : chenillard}

Le tableau~\ref{tab:p3_tests} donne, pour chaque flanc montant de
\texttt{clock\_i}, la fonction testée et la valeur attendue sur \texttt{LED\_out}.
Les valeurs observées sur le chronogramme (figures~\ref{fig:p3_chrono_a} et~\ref{fig:p3_chrono_b})
correspondent toutes aux valeurs attendues.

\begin{table}[H]
  \centering
  \small
  \begin{tabular}{c l c}
    \hline
    \textbf{Flanc} & \textbf{Fonction testée} & \textbf{\texttt{LED\_out} attendu} \\
    \hline
    5~µs   & Reset actif (\texttt{nReset\_i} = 0)               & \texttt{0000 0000} \\
    30~µs  & Chargement du pattern \texttt{10000000}            & \texttt{1000 0000} \\
    50~µs  & Maintien (\texttt{enable\_i} = 0)                  & \texttt{1000 0000} \\
    65 à 125~µs      & Rotation à droite (7 flancs)                       & \texttt{0100 0000} $\rightarrow$ \texttt{0000 0001} \\
    135~µs & \textbf{Rebouclage à droite} & \texttt{1000 0000} \\
    145~µs & \textbf{Rebouclage à gauche} & \texttt{0000 0001} \\
    155 et 165~µs     & Rotation à gauche                                  & \texttt{0000 0010}, \texttt{0000 0100} \\
    175~µs & Maintien (\texttt{enable\_i} = 0)                  & \texttt{0000 0100} \\
    185~µs & \textbf{Load prioritaire} (\texttt{11010000}) & \texttt{1101 0000} \\
    200 à 230~µs     & Rotation à droite            & \texttt{0110 1000} $\rightarrow$ \texttt{0000 1101} \\
    240~µs & Rebouclage avec plusieurs bits actifs              & \texttt{1000 0110} \\
    255~µs & Changement de sens (gauche)                        & \texttt{0000 1101} \\
    260~µs & \textbf{Reset asynchrone}      & \texttt{0000 0000} \\
    \hline
  \end{tabular}
  \caption{Tests de simulation du chenillard}
  \label{tab:p3_tests}
\end{table}

\capture{p3_chrono_a}{Chronogramme de \texttt{chenillard\_sim} (1/2) : de 0 à 138~µs}{fig:p3_chrono_a}{}
\capture{p3_chrono_b}{Chronogramme de \texttt{chenillard\_sim} (2/2) : de 132 à 266~µs}{fig:p3_chrono_b}{}
